\documentclass[12pt, a4paper]{article}

% ─── Packages ─────────────────────────────────────────────────────────────────
\usepackage[utf8]{inputenc}
\usepackage[T1]{fontenc}
\usepackage[french]{babel}
\usepackage{amsmath, amssymb}
\usepackage{geometry}
\usepackage{hyperref}
\usepackage{booktabs}
\usepackage{array}
\usepackage{longtable}
\usepackage{graphicx}
\usepackage{cite}
\usepackage{setspace}
\usepackage{parskip}
\usepackage{titlesec}
\usepackage{fancyhdr}
\usepackage{xcolor}

\geometry{top=2.5cm, bottom=2.5cm, left=2.5cm, right=2.5cm}
\onehalfspacing

\hypersetup{
    colorlinks=true,
    linkcolor=blue!60!black,
    citecolor=blue!60!black,
    urlcolor=blue!60!black
}

% ─── En-tête ─────────────────────────────────────────────────────────────────
\pagestyle{fancy}
\fancyhf{}
\rhead{\small Détection de variations atypiques par LOF  BEAC}
\lhead{\small Méthodologie}
\rfoot{\thepage}

% ─── Titre ───────────────────────────────────────────────────────────────────
\title{
    \textbf{Méthodologie Modèle 1 : Détection d'Anomalies par LOF\\
    sur les Données Monétaires BEAC}\\[0.3em]
}
\author{
    \textbf{MOUKOKO EKAMBI GEORGES MIGUEL}\\[0.3em]
    \textit{Stagiaire, Banque des États de l'Afrique Centrale}\\
    \textit{Données sources : Banque des États de l'Afrique Centrale (BEAC)}
}

\date{\today}

% ─── Document ────────────────────────────────────────────────────────────────
\begin{document}

\maketitle
\tableofcontents
\newpage

% ─────────────────────────────────────────────────────────────────────────────
\section{Architecture du Pipeline Multi-fichiers}
\label{sec:architecture}
% ─────────────────────────────────────────────────────────────────────────────

\subsection{Inventaire des sources de données}
\label{subsec:inventaire}

Le dossier de données contient douze fichiers au format XLSX, organisés selon
la nomenclature
\texttt{clean\_<pays>\_beac\_mapping\_2SR\_<Actif|Passif>.xlsx},
et couvrant l'ensemble des six États membres de la CEMAC
(tableau~\ref{tab:fichiers}).

\begin{table}[htbp]
    \centering
    \caption{Fichiers sources disponibles par pays et par volet comptable.}
    \label{tab:fichiers}
    \begin{tabular}{lcc}
        \toprule
        \textbf{Pays} & \textbf{Actif} & \textbf{Passif} \\
        \midrule
        Cameroun            & $\checkmark$ & $\checkmark$ \\
        Congo               & $\checkmark$ & $\checkmark$ \\
        Gabon               & $\checkmark$ & $\checkmark$ \\
        Guinée équatoriale  & $\checkmark$ & $\checkmark$ \\
        RCA                 & $\checkmark$ & $\checkmark$ \\
        Tchad               & $\checkmark$ & $\checkmark$ \\
        \bottomrule
    \end{tabular}
\end{table}

Chaque fichier partage la même structure formelle (format large, période
mensuelle \texttt{AAAMyy}, indicateurs IFS), mais peut présenter des
hétérogénéités de couverture temporelle et de taux de complétude selon
le pays et le volet comptable. Les pays à système bancaire moins développé
(RCA, Guinée équatoriale) présentent a priori des taux de valeurs manquantes
supérieurs et des fenêtres temporelles plus courtes. Le pipeline est conçu
pour s'adapter à ces variations sans intervention manuelle.

\subsection{Choix du mode d'exécution : pipeline séparé}
\label{subsec:mode_execution}

Trois architectures d'exécution du LOF sont envisageables pour traiter
l'ensemble de ces fichiers :

\begin{enumerate}
    \item \textbf{Mode séparé} : le pipeline complet (prétraitement,
    décomposition, réduction, LOF) est exécuté indépendamment sur chacun
    des douze fichiers. On obtient douze jeux de scores LOF, chacun mesurant
    l'anormalité d'un mois \textit{relativement à la dynamique propre de
    ce pays et de ce volet comptable}.

    \item \textbf{Mode pool CEMAC} : les douze fichiers sont concaténés en
    un unique jeu de données avant l'exécution du LOF. L'anormalité est
    alors mesurée relativement à l'ensemble des observations de tous les pays,
    sur les deux bilans simultanément.

    \item \textbf{Mode hiérarchique} : combinaison des deux niveaux
    précédents, avec agrégation des scores par pays puis synthèse au
    niveau CEMAC.
\end{enumerate}

Nous retenons le mode séparé pour trois raisons complémentaires.

\textit{Premièrement, la comparabilité inter-pays est structurellement
inadéquate.}
Les économies CEMAC présentent des structures bancaires très hétérogènes
(taille du secteur financier, degré d'intermédiation, profondeur des marchés) :
une même valeur absolue d'un agrégat peut être parfaitement normale au Cameroun
et extrême en RCA. Concaténer les fichiers sans normalisation inter-pays
préalable revient à mélanger des distributions incomparables dans un espace
euclidien commun, ce qui génère des scores LOF mesurant l'écart à la moyenne
CEMAC plutôt que l'anomalie propre à chaque pays.

\textit{Deuxièmement, la calibration de MinPts serait invalidée.}
Dans un pool CEMAC ($n \approx 195 \times 12 = 2\,340$), le paramètre
$\text{MinPts}$ devrait être recalibré sur une base de données douze fois
plus grande, et son interprétation en termes de \og{}cycles de
référence\fg{} (section~\ref{sec:calibration}) perdrait son ancrage
économique.

\textit{Troisièmement, seul le mode séparé produit une interprétation
économique directe.}
Un score LOF élevé en avril~2020 au Cameroun doit être lu comme une anomalie
dans la dynamique bancaire camerounaise, et non comme une position marginale
dans un espace CEMAC composite qui mêle des régimes économiques sans lien
causal.

Le mode pool CEMAC est néanmoins conservé comme variante d'exploration pour
la détection d'anomalies \textit{systémiques}, c'est-à-dire d'événements
simultanément anormaux dans plusieurs pays
(section~\ref{subsec:protocole_val}).

\subsection{Cohérence comptable Actif--Passif}
\label{subsec:actif_passif}

L'identité fondamentale du bilan, $\text{Total Actif} = \text{Total Passif}$,
constitue une contrainte comptable exploitable comme mécanisme de validation
interne. Pour chaque pays, les scores LOF obtenus sur le fichier Actif seront
confrontés à ceux obtenus sur le fichier Passif.

Deux situations sont distinguées :
\begin{itemize}
    \item \textbf{Anomalie concordante} (signal sur Actif \textit{et} Passif
    au même mois) : la déviation affecte la structure du bilan dans son
    ensemble  c'est le signe d'un événement économique réel (choc de
    crédit, assèchement de liquidités, crise systémique).
    \item \textbf{Anomalie discordante} (signal sur un seul volet) :
    suggère prioritairement une erreur de collecte, de saisie ou un problème
    de réconciliation comptable entre les deux fichiers, plutôt qu'un
    événement économique réel.
\end{itemize}

Cette distinction est intégrée au protocole de validation de la
section~\ref{sec:validation}.

% ─────────────────────────────────────────────────────────────────────────────
\section{Prétraitement et Mise en Forme des Données Monétaires}
\label{sec:pretraitement}
% ─────────────────────────────────────────────────────────────────────────────

\subsection{Structure brute du dataset}
\label{subsec:structure_brute}

Les douze fichiers sources sont organisés en format large
(\textit{wide format}) : chaque ligne représente un indicateur financier
identifié par son code IFS, et chaque colonne correspond à une période
mensuelle notée au format \texttt{AAAMyy}
(ex.\ \texttt{2010M1}, \texttt{2024M12}).
À l'inspection du fichier de référence (Cameroun Actif), la matrice de données
couvre la période de décembre~2001 à mars~2026, soit environ $195$
observations temporelles, pour $\approx 55$ indicateurs financiers répartis
en grandes catégories d'actif : Numéraire et Dépôts, Titres, Crédits, Actions,
Réserves Techniques d'Assurance, Autres Comptes à Recevoir et Crédits
Commerciaux. Les valeurs sont exprimées en milliards de FCFA
(Franc CFA de la CEMAC). Ces dimensions sont indicatives ; les autres fichiers
peuvent présenter des couvertures temporelles et des nombres d'indicateurs
différents, selon le pays et le volet comptable.

\subsection{Anomalies structurelles et valeurs manquantes}
\label{subsec:anomalies_struct}

L'examen des fichiers sources révèle deux types de problèmes structurels
distincts qu'il convient de traiter avant toute modélisation :

\begin{enumerate}
    \item \textbf{Valeurs \texttt{\#VALUE!}} : certaines cellules contiennent
    des erreurs de calcul héritées des formules Excel sources
    (ex.\ séries \textit{Numéraire et Dépôts}, \textit{Dépôts Transférables},
    \textit{Crédits} sur les périodes \texttt{2012M7}--\texttt{2012M9} et
    ponctuellement sur d'autres fenêtres). Ces entrées doivent être
    considérées comme des valeurs manquantes (\texttt{NaN}) et traitées
    par imputation.

    \item \textbf{Cellules vides} : plusieurs indicateurs de sous-catégorie
    présentent une première observation nulle ou absente pour la période
    \texttt{2001M12}, correspondant à une couverture temporelle incomplète
    en début de série.
\end{enumerate}

\subsection{Protocole de nettoyage}
\label{subsec:protocole}

La matrice des données est pivotée en format long (\textit{long format}),
où chaque ligne représente un triplet (indicateur, période, valeur),
facilitant les opérations d'imputation et de normalisation.
Le protocole de nettoyage comporte deux étapes.

\paragraph{Imputation des valeurs manquantes.}

La stratégie d'imputation doit être adaptée au \textit{mécanisme de
génération} des données manquantes, dont la nature conditionne la validité
statistique des méthodes disponibles.

\begin{itemize}
    \item \textbf{Lacunes courtes internes (1 à 3 mois consécutifs).}
    Ces lacunes, imputables aux erreurs de calcul Excel ou à des retards
    de remontée, peuvent raisonnablement être supposées
    \textit{Missing At Random} (MAR). L'\textbf{interpolation temporelle
    linéaire} est appliquée : elle préserve la continuité locale de la
    série sans introduire de biais structurel, en supposant une évolution
    approximativement linéaire entre deux observations adjacentes connues.

    \item \textbf{Lacunes étendues en milieu de série (plus de 3 mois
    consécutifs).}
    L'\textbf{algorithme MICE} (\textit{Multivariate Imputation by Chained
    Equations}) est mobilisé \cite{vanbuuren2011mice}. Il exploite les
    corrélations inter-séries pour inférer les valeurs absentes à partir de
    l'information contenue dans les autres indicateurs du bilan, en itérant
    des régressions conditionnelles sur chaque variable manquante. La validité
    de MICE repose sur l'hypothèse MAR, raisonnable ici : les lacunes
    résultent de défaillances de collecte liées à la période ou à
    l'indicateur (observables), et non à la valeur non observée elle-même.

    \item \textbf{Lacunes structurelles en début ou fin de série.}
    Ces lacunes correspondent à des périodes antérieures au démarrage
    effectif de la collecte d'un indicateur donné, phénomène fréquent en
    RCA et en Guinée équatoriale. Elles sont de nature
    \textit{Missing Not At Random} (MNAR). Appliquer MICE sur des données
    MNAR introduit un biais d'imputation non contrôlé et potentiellement
    sévère.
    \textbf{Traitement retenu} : troncature de la fenêtre temporelle
    d'analyse au premier mois effectivement renseigné pour la majorité des
    indicateurs du fichier considéré. Les séries dont la part de valeurs
    imputées en début de période dépasse $20\,\%$ de la fenêtre totale sont
    exclues de la matrice d'entrée du LOF pour ce fichier.
\end{itemize}

\paragraph{Standardisation robuste par série (médiane\,/\,MAD).}

Les $\approx 55$ indicateurs étant exprimés en milliards de FCFA mais
présentant des ordres de grandeur très hétérogènes, une standardisation
individuelle est impérative avant tout calcul de distance.

\textbf{Pourquoi le Z-score classique est inadapté ici.}
La standardisation usuelle $x^* = (x - \bar{x}) / \sigma$ utilise la moyenne
$\bar{x}$ et l'écart-type $\sigma$, tous deux très sensibles aux valeurs
extrêmes : un outlier prononcé gonfle simultanément $\bar{x}$ et $\sigma$,
réduisant mécaniquement sa propre valeur normalisée. Il devient donc moins
détectable après standardisation  précisément après l'étape censée préparer
sa détection. Leys \textit{et al.}~\cite{leys2013mad} formalisent cette
circularité et montrent sur données réelles que le Z-score peut masquer des
anomalies que la standardisation robuste révèle.

\textbf{Standardisation retenue : médiane\,/\,MAD.}
Pour chaque série $i$, la transformation appliquée est :
\[
    x_{i,t}^* = \frac{x_{i,t} - \tilde{x}_i}
                     {\kappa \cdot \operatorname{MAD}(x_i)}
\]
où $\tilde{x}_i = \operatorname{median}_t(x_{i,t})$ est la médiane
temporelle,
\[
    \operatorname{MAD}(x_i) =
    \operatorname{median}_t\!\bigl(\,|x_{i,t} - \tilde{x}_i|\,\bigr)
\]
est la déviation absolue médiane, et
$\kappa = 1/\Phi^{-1}(0{,}75) \approx 1{,}4826$ est le facteur de cohérence
rendant le MAD asymptotiquement équivalent à $\sigma$ sous distribution
gaussienne \cite{leys2013mad}. La médiane possède un point de rupture de
$50\,\%$ contre $0\,\%$ pour la moyenne, et le MAD hérite de cette propriété.

% ─────────────────────────────────────────────────────────────────────────────
\section{Décomposition et Stationnarisation des Séries Temporelles}
\label{sec:decomposition}
% ─────────────────────────────────────────────────────────────────────────────

\subsection{Inadéquation du LOF sur séries brutes}
\label{subsec:inadequation_lof}

Le Local Outlier Factor, tel que formalisé par Breunig \textit{et al.}
\cite{breunig2000lof}, opère dans un espace géométrique statique : il évalue
la densité locale d'un point en fonction de ses distances aux $k$ plus proches
voisins dans l'espace des descripteurs. Or, les agrégats monétaires du bilan
2SR de la BEAC exhibent des tendances macroéconomiques prononcées sur la
période 2001--2026 : par exemple, les crédits aux Autres Sociétés Non
Financières ont progressé de près de $600\,\%$ en termes nominaux. En présence
de telles tendances, la distance géométrique entre une observation de 2010 et
une observation de 2024 intègre mécaniquement l'effet de la dérive temporelle,
et non l'anomalie conjoncturelle recherchée.

La même problématique s'applique à la saisonnalité : si les dépôts
transférables présentent des pics récurrents en décembre (clôtures budgétaires)
et des creux en janvier, un LOF appliqué aux données brutes qualifiera ces
fluctuations saisonnières normales d'anomalies, faute de pouvoir les distinguer
de déviations genuinement atypiques.

\subsection{Décomposition STL}
\label{subsec:stl}

Pour isoler la composante résiduelle stationnaire sur laquelle le LOF sera
exécuté, nous mobilisons la décomposition STL (\textit{Seasonal and Trend
decomposition using Loess}), introduite par Cleveland \textit{et al.}
\cite{cleveland1990stl}, qui décompose chaque série $x_{i,t}$ en trois
composantes additives :
\[
    x_{i,t} = T_{i,t} + S_{i,t} + R_{i,t}
\]
où $T_{i,t}$ est la tendance estimée par régression locale (Loess),
$S_{i,t}$ est la composante saisonnière mensuelle (période $= 12$),
et $R_{i,t}$ est le résidu.

Deux alternatives classiques auraient pu être retenues : la différenciation
première et la méthode X-13~ARIMA-SEATS (standard des banques centrales).
STL est préféré pour trois raisons : il s'adapte aux tendances non linéaires ;
sa composante saisonnière est \textit{time-varying}, pertinente pour des séries
dont la saisonnalité budgétaire peut évoluer avec les réformes fiscales CEMAC ;
enfin, il est robuste aux outliers isolés via un schéma de pondération interne,
ce qui évite que les anomalies recherchées biaisent l'estimation de la
tendance.

Seul le résidu $R_{i,t}$ est retenu comme entrée du LOF. Ce résidu représente
les fluctuations inexpliquées par la dynamique normale de la série, et
constitue donc le support naturel de la détection d'anomalies, conformément
aux recommandations de l'état de l'art du projet \cite{moukoko2025etatart}.

% ─────────────────────────────────────────────────────────────────────────────
\section{Réduction de la Dimensionnalité}
\label{sec:reduction}
% ─────────────────────────────────────────────────────────────────────────────

\subsection{Le fléau de la dimensionnalité}
\label{subsec:fleau_dim}

Avec $p \approx 55$ indicateurs financiers, l'espace résiduel dans lequel
le LOF calcule ses distances est de dimension~55. Dans un espace de haute
dimension, la distance euclidienne entre deux points tend vers la même valeur
pour l'ensemble des paires de points (phénomène de concentration des
distances), rendant le concept de \og{}plus proche voisin\fg{} de moins en
moins discriminant à mesure que $p$ croît \cite{breunig2000lof}, Section~7.4.
Ce phénomène est analysé de façon systématique par Zimek \textit{et al.}
\cite{zimek2012survey}, qui montrent que la plupart des méthodes à base de
densité locale, dont le LOF, dégénèrent en dimension supérieure à~20 sans
réduction préalable.

Dans le contexte du dataset BEAC, ce problème est amplifié par la redondance
structurelle entre indicateurs : de nombreuses séries sont des composantes ou
des sous-totaux d'agrégats de niveau supérieur (ex.\ \textit{Dépôts
Transférables en MN} et ses sous-postes par contrepartie), induisant une
forte collinéarité qui concentre la variance dans un sous-espace de faible
dimension.

\subsection{Stratégie de projection}
\label{subsec:projection}

\begin{enumerate}
    \item \textbf{ACP Robuste (Robust PCA  RPCA).}
    L'ACP classique est sensible aux outliers : ceux-ci exercent un effet de
    levier sur les vecteurs propres, tirant les axes principaux vers eux et
    biaisent la projection, si bien que les anomalies se retrouvent projetées
    dans le sous-espace jugé \og{}normal\fg{}, réduisant leur visibilité dans
    l'espace de calcul du LOF. Nous privilégions donc la RPCA, formalisée par
    Candès \textit{et al.}~\cite{candes2011rpca}, qui décompose la matrice
    $X \in \mathbb{R}^{n \times p}$ en une composante de bas rang $L$
    (structure normale) et une composante creuse $S$ (anomalies) :
    \[
        X = L + S
    \]
    par la minimisation convexe :
    \[
        \min_{L,\,S} \; \|L\|_* + \lambda\|S\|_1
        \quad \text{s.c.} \quad L + S = X
    \]
    où $\|\cdot\|_*$ est la norme nucléaire et $\|\cdot\|_1$ la norme
    \textit{entrywise} $\ell_1$. Seule la composante $L$ est ensuite
    projetée par ACP classique pour obtenir les $k$ premières composantes
    principales, $k$ étant déterminé par le critère du coude et validé par
    le pourcentage de variance expliquée cumulée (cible : $90\,\%$).

    \item \textbf{UMAP} (\textit{Uniform Manifold Approximation and
    Projection}) : méthode non linéaire réservée à la \textit{visualisation}
    des scores LOF. Son usage est limité à l'exploration graphique : sa nature
    non déterministe et la distorsion qu'il introduit dans les distances
    globales le rendent inadapté comme espace de calcul du LOF dans un
    contexte opérationnel où la reproductibilité est exigée.
\end{enumerate}

% ─────────────────────────────────────────────────────────────────────────────
\section{Formalisation de la Densité Locale et du Facteur LOF}
\label{sec:lof_formel}
% ─────────────────────────────────────────────────────────────────────────────

Cette section présente, dans leur enchaînement conceptuel formel,
les définitions introduites par Breunig \textit{et al.}~\cite{breunig2000lof}
qui fondent le calcul du facteur LOF.

\subsection{\textit{k}-distance et voisinage}
\label{subsec:kdistance}

Pour un entier $k \geq 1$ et un point $p$ de l'espace des résidus normalisés,
la \textbf{\textit{k}-distance} de $p$, notée $k\text{-distance}(p)$, est
la distance $d(p, o)$ vers un objet $o$ tel qu'au moins $k$ objets
$o' \neq p$ vérifient $d(p, o') \leq d(p, o)$ et au plus $k-1$ vérifient
$d(p, o') < d(p, o)$ \cite{breunig2000lof}, Déf.~3.

Le \textbf{voisinage de \textit{k}-distance} de $p$ est :
\[
    \mathcal{N}_k(p) = \bigl\{\, q \in \mathcal{D} \setminus \{p\}
    \mid d(p, q) \leq k\text{-distance}(p) \,\bigr\}
\]
\cite{breunig2000lof}, Déf.~4. La cardinalité de $\mathcal{N}_k(p)$ peut
être supérieure à $k$ en cas d'équidistance de plusieurs points.

\subsection{Distance de joignabilité}
\label{subsec:reachability}

La \textbf{distance de joignabilité} de $p$ par rapport à $o$
\cite{breunig2000lof}, Déf.~5, est :
\[
    \text{reach-dist}_k(p, o) = \max \bigl\{
    k\text{-distance}(o),\; d(p, o)
    \bigr\}
\]
Si $p$ est éloigné de $o$, la distance réelle $d(p,o)$ est retenue ;
si $p$ est très proche de $o$, la distance est plafonnée à
$k\text{-distance}(o)$, réduisant les fluctuations statistiques locales.

\subsection{Densité locale de joignabilité}
\label{subsec:lrd}

La \textbf{densité locale de joignabilité} de $p$
\cite{breunig2000lof}, Déf.~6, est :
\[
    \text{lrd}_k(p) = \left(
    \frac{\displaystyle\sum_{o \in \mathcal{N}_k(p)}
    \text{reach-dist}_k(p, o)}
    {\left|\mathcal{N}_k(p)\right|}
    \right)^{-1}
\]
Un $\text{lrd}_k(p)$ élevé indique que $p$ est situé dans une région dense,
c'est-à-dire que son comportement mensuel est similaire à celui de plusieurs
périodes voisines  un comportement normal dans le contexte BEAC.

\subsection{Facteur d'anomalie local (LOF)}
\label{subsec:lof_def}

Le \textbf{facteur d'anomalie local} de $p$
\cite{breunig2000lof}, Déf.~7, est :
\[
    \text{LOF}_k(p) =
    \frac{\displaystyle\sum_{o \in \mathcal{N}_k(p)}
    \dfrac{\text{lrd}_k(o)}{\text{lrd}_k(p)}}
    {\left|\mathcal{N}_k(p)\right|}
\]
Un $\text{LOF}_k(p) \approx 1$ caractérise un point normal ; un
$\text{LOF}_k(p) \gg 1$ signale une anomalie locale. Breunig \textit{et al.}
établissent des bornes supérieure et inférieure générales sur ce facteur en
fonction des distances de joignabilité directes et indirectes
\cite{breunig2000lof}, Théorèmes~1 et~2. Appliqué au dataset BEAC, un mois
$t$ recevra un $\text{LOF}_k$ élevé si son résidu est isolé dans l'espace
réduit par RPCA, c'est-à-dire si ce mois constitue une déviation simultanée
et atypique de plusieurs séries financières par rapport à leur comportement
habituel.

% ─────────────────────────────────────────────────────────────────────────────
\section{Calibration Empirique du Paramètre \textit{MinPts}}
\label{sec:calibration}
% ─────────────────────────────────────────────────────────────────────────────

\subsection{Non-monotonicité du LOF en fonction de \textit{MinPts}}
\label{subsec:non_monotonicite}

Le seul hyperparamètre du LOF est $\text{MinPts}$, qui fixe le nombre de
voisins utilisé pour estimer la densité locale. Breunig \textit{et al.}
\cite{breunig2000lof}, Section~6.1, démontrent que le score LOF n'évolue pas
de manière monotone avec $\text{MinPts}$ : pour un même point, augmenter
$\text{MinPts}$ peut alternativement augmenter, diminuer ou stabiliser son
LOF, selon la structure locale du voisinage (Figure~7 de l'article de
référence).

\subsection{Stratégie de la plage de valeurs}
\label{subsec:plage_valeurs}

Pour pallier cette instabilité, Breunig \textit{et al.}
\cite{breunig2000lof}, Section~6.2, proposent d'utiliser une plage de valeurs
$[\text{MinPts}_\text{LB},\, \text{MinPts}_\text{UB}]$ et de retenir pour
chaque point le score maximum :
\[
    \text{LOF}^*(p) = \max \bigl\{
    \text{LOF}_{\text{MinPts}}(p) \mid
    \text{MinPts}_\text{LB} \leq \text{MinPts} \leq \text{MinPts}_\text{UB}
    \bigr\}
\]
Le maximum met en évidence l'instance à laquelle le point est le plus isolé,
tandis que la moyenne diluerait ce caractère atypique et le minimum pourrait
l'effacer complètement \cite{breunig2000lof}.

\subsection{Détermination de la borne inférieure
\textit{MinPts}$_\text{LB}$}
\label{subsec:borne_inf}

\begin{enumerate}
    \item \textbf{Critère de stabilité statistique} : la règle pratique de
    Breunig \textit{et al.}~\cite{breunig2000lof} est de fixer
    $\text{MinPts}_\text{LB} \geq 10$ pour éliminer les fluctuations
    statistiques indésirables. Avec $n \approx 195$, la valeur~10 représente
    le 5\textsuperscript{e} percentile de la distribution des tailles de
    voisinage.

    \item \textbf{Critère de taille minimale de cluster} :
    $\text{MinPts}_\text{LB}$ correspond au nombre minimal de mois
    définissant un régime normal de référence. Une valeur de~12 est envisagée,
    correspondant à un cycle annuel complet, cohérente avec la saisonnalité
    budgétaire de la zone CEMAC.
\end{enumerate}

\subsection{Détermination de la borne supérieure
\textit{MinPts}$_\text{UB}$}
\label{subsec:borne_sup}

$\text{MinPts}_\text{UB}$ représente la taille maximale d'un groupe de points
pouvant être considéré comme un cluster d'anomalies locales
\cite{breunig2000lof}. Nous proposons $\text{MinPts}_\text{UB} = 30$
à~$40$, correspondant à 2 à 3 ans de données, motivés par la durée typique
des cycles macroéconomiques CEMAC. La plage $[10, 30]$ sera validée par
analyse de sensibilité.

Pour tout fichier avec $n < 100$ observations effectives (RCA, Guinée
équatoriale), la borne supérieure est ramenée à
$\text{MinPts}_\text{UB} = \lfloor n/5 \rfloor$ afin de maintenir un rapport
$\text{MinPts}/n$ cohérent avec les recommandations de l'article de référence.

% ─────────────────────────────────────────────────────────────────────────────
\section{Protocole de Validation et de Comparaison}
\label{sec:validation}
% ─────────────────────────────────────────────────────────────────────────────

\subsection{Contexte non supervisé : absence de labels}
\label{subsec:contexte_ns}

La détection d'anomalies dans le dataset BEAC est intrinsèquement non
supervisée : il n'existe pas de labels historiques certifiés. Breunig
\textit{et al.}~\cite{breunig2000lof}, Section~7, soulignent cette dimension
exploratoire du LOF : il s'agit d'un outil de \textit{scoring} ordinal dont
la valeur réside dans le classement relatif des observations, et non dans un
seuil absolu.

\subsection{Formalisation du seuil de détection $\tau$}
\label{subsec:seuil}

Le seuil $\tau$ est ancré dans la distribution empirique des scores
LOF$^*$ \cite{kriegel2011outlier} via la règle de l'écart interquartile :
\[
    \tau = Q_3\!\bigl(\text{LOF}^*\bigr)
    + 1{,}5 \times \operatorname{IQR}\!\bigl(\text{LOF}^*\bigr)
\]
où $Q_3$ est le troisième quartile et
$\operatorname{IQR} = Q_3 - Q_1$ \cite{tukey1977eda}.
Ce facteur est préféré au percentile fixe~95, dont la sensibilité dépend de
la proportion d'anomalies réelles dans le dataset  inconnue a priori. En
complément, le seuil au 95\textsuperscript{e} percentile est calculé comme
variante de sensibilité. Les observations dépassant les deux seuils
simultanément constituent les \textit{anomalies robustes} ; celles signalées
par un seul seuil constituent les \textit{anomalies marginales} requérant
une expertise domaine.

\subsection{Protocole de validation proposé}
\label{subsec:protocole_val}

En l'absence de labels, le protocole de validation repose sur cinq niveaux
complémentaires :

\begin{enumerate}
    \item \textbf{Analyse de sensibilité des scores LOF.}
    Le profil $\text{LOF}_{\text{MinPts}}(t)$ est tracé pour
    $\text{MinPts} \in [10, 30]$, en reproduisant la démarche de la
    Figure~8 de \cite{breunig2000lof} sur les données BEAC. Les points
    maintenant un LOF élevé sur toute la plage sont des anomalies robustes.

    \item \textbf{Confrontation aux ruptures macroéconomiques.}
    Les anomalies détectées ($\text{LOF}^* > \tau$, seuil défini à la
    section~\ref{subsec:seuil}) sont confrontées aux événements documentés
    dans la zone CEMAC : choc pétrolier (2014--2016), crise sécuritaire au
    Cameroun (à partir de 2017), choc COVID-19 (2020--2021), dynamiques
    post-pandémiques (2022--2024). Une concordance entre les anomalies
    détectées et ces épisodes constitue une forme de validation externe
    qualitative.

    \item \textbf{Comparaison avec une baseline de distance globale.}
    Le LOF est comparé à l'approche \textit{DB(pct, dmin)-outlier}
    au sens de Knorr et Ng \cite{breunig2000lof}, Section~3, appliquée aux
    mêmes résidus, afin de mettre en évidence les anomalies locales que le
    LOF détecte et que l'approche globale manque, notamment les déviations
    atypiques dans des régimes à densité variable
    (ex.\ post-2020 vs pré-2015).

    \item \textbf{Validation croisée inter-pays.}
    Les douze jeux de scores LOF$^*$ sont superposés sur l'axe temporel.
    Un mois $t$ est qualifié d'\textbf{anomalie systémique CEMAC} s'il
    dépasse $\tau$ dans au moins $m \geq 3$ pays simultanément. Cette
    procédure exploite une opportunité unique offerte par la structure
    multi-fichiers et constitue une validation externe quantitative : les
    anomalies systémiques devraient coïncider avec les chocs documentés
    (crise pétrolière, COVID-19). À l'inverse, une anomalie détectée dans
    un seul pays sur une période non documentée est plus vraisemblablement
    liée à un problème de collecte local qu'à un événement économique réel.

    \item \textbf{Cohérence Actif--Passif.}
    Les scores LOF$^*$ des fichiers Actif et Passif sont comparés mois par
    mois pour chaque pays. Les anomalies concordantes (signal sur les deux
    volets) sont interprétées comme des événements économiques réels ; les
    anomalies discordantes sont signalées comme suspectes et soumises à
    vérification comptable (section~\ref{subsec:actif_passif}).
\end{enumerate}

\subsection{Présentation des résultats}
\label{subsec:presentation}

Les résultats seront présentés sous la forme d'une courbe de scores
$\text{LOF}^*(t)$ en fonction du temps, superposée à la série brute de
l'indicateur considéré, permettant d'associer visuellement chaque anomalie
détectée à son contexte chronologique et économique. Un tableau récapitulatif
identifiera les $N$ périodes les plus anomales (classées par $\text{LOF}^*$
décroissant), avec pour chacune le ou les indicateurs du bilan portant la
déviation principale, une interprétation économique proposée, et la mention
de sa concordance inter-pays et inter-bilans.

% ═════════════════════════════════════════════════════════════════════════════
% NOUVELLES RÉFÉRENCES BIBLIOGRAPHIQUES
%
% Ces entrées sont à FUSIONNER avec le \begin{thebibliography}{99} du document
% maître (ou à ajouter dans le fichier .bib si la compilation utilise BibTeX).
% Elles correspondent aux sept sources mobilisées dans ce fichier qui ne
% figurent pas encore dans la bibliographie de l'état de l'art.
% ═════════════════════════════════════════════════════════════════════════════

\bibliographystyle{plain}
\begin{thebibliography}{99}

% ─── Référence principale (LOF) ───────────────────────────────────────────────
\bibitem{breunig2000lof}
Breunig, M.M., Kriegel, H.-P., Ng, R.T., \& Sander, J. (2000).
\textit{LOF: Identifying Density-Based Local Outliers}.
Proceedings of the 2000 ACM SIGMOD International Conference on
Management of Data, pp.~93--104.
\url{https://doi.org/10.1145/342009.335388}

% ─── Document compagnon du projet ─────────────────────────────────────────────
\bibitem{moukoko2025etatart}
Moukoko~Ekambi, G.M. (2025).
\textit{Détection de Variations Atypiques dans les Séries Temporelles
Financières par l'Intelligence Artificielle : État de l'Art}.
Rapport de stage, Banque des États de l'Afrique Centrale (BEAC).
ENSPD / Université de Douala, LIDIA Laboratory.

% ─── Références méthodologiques nouvelles ─────────────────────────────────────
\bibitem{cleveland1990stl}
Cleveland, R.B., Cleveland, W.S., McRae, J.E., \& Terpenning, I. (1990).
\textit{STL: A Seasonal-Trend Decomposition Procedure Based on Loess}.
Journal of Official Statistics, \textbf{6}(1), 3--73.

\bibitem{vanbuuren2011mice}
Van~Buuren, S., \& Groothuis-Oudshoorn, K. (2011).
\textit{mice: Multivariate Imputation by Chained Equations in R}.
Journal of Statistical Software, \textbf{45}(3), 1--67.
\url{https://doi.org/10.18637/jss.v045.i03}

\bibitem{leys2013mad}
Leys, C., Ley, C., Klein, O., Bernard, P., \& Licata, L. (2013).
\textit{Detecting outliers: Do not use standard deviation around the mean,
use absolute deviation around the median}.
Journal of Experimental Social Psychology, \textbf{49}(4), 764--766.
\url{https://doi.org/10.1016/j.jesp.2013.03.013}

\bibitem{zimek2012survey}
Zimek, A., Schubert, E., \& Kriegel, H.-P. (2012).
\textit{A survey on unsupervised outlier detection in high-dimensional
numerical data}.
Statistical Analysis and Data Mining, \textbf{5}(5), 363--387.
\url{https://doi.org/10.1002/sam.11161}

\bibitem{candes2011rpca}
Cand\`{e}s, E.J., Li, X., Ma, Y., \& Wright, J. (2011).
\textit{Robust Principal Component Analysis?}
Journal of the ACM, \textbf{58}(3), Article~11.
\url{https://doi.org/10.1145/1970392.1970395}

\bibitem{kriegel2011outlier}
Kriegel, H.-P., Kr\"{o}ger, P., Schubert, E., \& Zimek, A. (2011).
\textit{Interpreting and Unifying Outlier Scores}.
Proceedings of the 2011 SIAM International Conference on Data Mining
(SDM 2011), pp.~13--24.
\url{https://doi.org/10.1137/1.9781611972818.2}

\bibitem{tukey1977eda}
Tukey, J.W. (1977).
\textit{Exploratory Data Analysis}.
Addison-Wesley, Reading (MA).

\end{thebibliography}

\end{document}